\documentclass[10pt,twocolumn]{article}

%  Packages 
\usepackage[margin=1in]{geometry}
\usepackage{times}
\usepackage{graphicx}
\usepackage{amsmath}
\usepackage{amssymb}
\usepackage{booktabs}
\usepackage{multirow}
\usepackage{xcolor}
\usepackage{hyperref}
\usepackage{listings}
\usepackage{caption}
\usepackage{subcaption}
\usepackage{microtype}
\usepackage{balance}
\usepackage{url}
\usepackage{cite}
\usepackage{array}
\usepackage{tabularx}

\hypersetup{
    colorlinks=true,
    linkcolor=blue!70!black,
    citecolor=blue!70!black,
    urlcolor=blue!70!black,
}

\lstset{
    basicstyle=\ttfamily\footnotesize,
    breaklines=true,
    frame=single,
    backgroundcolor=\color{gray!8},
    rulecolor=\color{gray!40},
    xleftmargin=4pt,
    xrightmargin=4pt,
}

\definecolor{resultgreen}{RGB}{21,128,61}
\definecolor{resultorange}{RGB}{180,90,0}

%  Title 
\title{%
    \textbf{HuddleCluster: A Penguin-Inspired Self-Organizing}\\
    \textbf{Load Balancer with Adaptive Thermal Eviction}%
}

\author{%
    Rahad Bhuiya\\
    Independent Researcher\\
    \texttt{github.com/rahadbhuiya/HuddleCluster}%
}

\date{}

%  Document 
\begin{document}

\maketitle

%  Abstract 
\begin{abstract}
We present \textsc{HuddleCluster}, a load-balancing algorithm inspired by the
thermoregulatory huddle behaviour of Emperor Penguins
(\textit{Aptenodytes forsteri}).
Traditional algorithms---Round Robin~(RR) and Least Connections~(LC)---lack
mechanisms for server rest periods and autonomous anomaly detection.
\textsc{HuddleCluster} organises servers into an active \emph{inner ring} and
a resting \emph{outer ring}, rotating servers based on a composite
Exponential Moving Average~(EMA) temperature score dominated by
\emph{relative latency anomaly}: how much slower a server is than the
cluster-wide median.

We evaluate the algorithm across simulated stress scenarios
(10~independent trials, Welch's $t$-test, $p{<}0.05$) and a
containerised HTTP deployment (6~FastAPI upstream servers).
Under server failure, \textsc{HuddleCluster} reduces P95 latency from
500\,ms to $23.9 \pm 0.5$\,ms in simulation and from 5\,026\,ms to
85.6\,ms in the real HTTP benchmark---both statistically significant
($p{<}0.001$).
Under normal load and traffic spikes, performance is comparable to RR
and LC\@.
Inner-ring fairness (Gini coefficient $= 0.00$) is achieved in every
scenario.
The routing overhead is 10.7\,$\mu$s per request with a 28.3\,KB peak
memory footprint for 20~servers.
Slow-server detection converges in 36~cluster requests on average
(range 35--40, 5~trials).
A containerised benchmark against NGINX (Round Robin) and NGINX (Least
Connections) shows P95 improvements of up to 39.4\% under normal load
and 35.3\% during server failure, confirming results against
production-grade baselines.
The implementation is available as open-source Python with zero runtime
dependencies at
\url{https://github.com/rahadbhuiya/HuddleCluster}.

\textbf{Keywords:} load balancing, distributed systems,
biologically-inspired computing, adaptive scheduling, anomaly detection,
fairness.
\end{abstract}

%  I. Introduction 
\section{Introduction}
\label{sec:intro}

Load balancing is a foundational problem in distributed computing.
As modern applications scale to millions of concurrent requests, the
choice of routing algorithm directly affects tail latency, resource
utilisation, and system resilience.
Despite decades of research, the most widely deployed algorithms---Round
Robin~(RR) and Least Connections~(LC)---share a fundamental limitation:
they are stateless with respect to server health over time.
Neither algorithm provides a mechanism for intentional server rest
periods, autonomous anomaly detection, or fairness guarantees.

This paper introduces \textsc{HuddleCluster}, a load-balancing data
structure inspired by the thermoregulatory behaviour of Emperor Penguins.
In Antarctic winters, penguins survive by forming dense huddles: cold
outer-edge penguins push inward, while warm inner penguins rotate
outward to rest---a self-organizing, leaderless process driven by local
temperature thresholds~\cite{zitterbart2011,waters2012}.

\textsc{HuddleCluster} maps this biological metaphor directly to server
scheduling.
Servers in the \emph{inner ring} actively serve requests; those in the
\emph{outer ring} rest and recover.
A server's \emph{temperature}---a composite EMA score derived from
relative latency anomaly, CPU, memory, connections, and error
rate---determines when it rotates out to rest and when it returns.

Our primary contributions are:
\begin{enumerate}
\item The \emph{dual-ring architecture} with temperature-driven rotation
      as a novel load-balancing primitive.
\item \emph{Relative latency anomaly scoring}, normalising observed
      latency against the cluster-wide median, making eviction decisions
      scale-invariant.
\item A complete open-source Python implementation with 45~passing
      unit tests and zero runtime dependencies.
\item Empirical evaluation across simulated and real HTTP scenarios with
      statistical significance testing.
\end{enumerate}

%  II. Background 
\section{Background and Related Work}
\label{sec:background}

\subsection{Classical Load-Balancing Algorithms}

\textbf{Round Robin} distributes requests cyclically, providing simple
fairness in request count but ignoring server state~\cite{eager1986}.
Under heterogeneous load, RR continues routing to degraded servers,
inflating tail latency.

\textbf{Least Connections} routes to the server with the fewest active
connections.
More adaptive than RR, but connection count does not capture latency or
CPU state: a server with few connections but high per-request latency is
indistinguishable from a fast idle server.

\textbf{Power-of-Two Choices}~\cite{mitzenmacher2001} selects two servers
at random and routes to the less loaded.
It improves over RR in expectation but provides no rest guarantees or
fairness tracking.

\subsection{Biologically-Inspired Scheduling}

Bio-inspired algorithms have a long history in distributed systems: Ant
Colony Optimisation for routing~\cite{dorigo1997}, flocking algorithms
for peer-to-peer topology management~\cite{reynolds1987}.
Direct application of penguin huddle dynamics to server scheduling has
not appeared in the literature to our knowledge.
The closest analogues are token-based work-stealing schedulers~\cite{blumofe1999},
which share the intuition of load migration but lack the ring structure
and temperature metric of \textsc{HuddleCluster}.

\subsection{Adaptive Load Balancing}

Production systems such as Maglev~\cite{eisenbud2016} and
PREQUAL~\cite{suresh2023} incorporate latency-aware routing and
connection-level health signals.
These systems require significant infrastructure investment.
\textsc{HuddleCluster} targets a different design point: a lightweight,
zero-dependency algorithm that self-adapts using only observable request
latency via a \texttt{record\_latency()} feedback API.

%  III. Algorithm Design 
\section{Algorithm Design}
\label{sec:design}

\subsection{Dual-Ring Architecture}

\textsc{HuddleCluster} partitions the server pool $S$ into two disjoint
sets at all times:
\begin{itemize}
\item \textbf{Inner ring} $I$: active servers handling requests,
      implemented as a \texttt{deque} for $O(1)$ circular round-robin.
\item \textbf{Outer ring} $O$: resting servers, implemented as a
      min-heap keyed by temperature for $O(\log n)$ promotion.
\end{itemize}
Ring sizes are bounded by \texttt{min\_inner} and \texttt{max\_inner}.
Figure~\ref{fig:architecture} illustrates the architecture.

\begin{figure}[t]
  \centering
  \includegraphics[width=\columnwidth]{architecture_diagram.png}
  \caption{%
    \textsc{HuddleCluster} dual-ring architecture.
    Inner ring servers (green) rotate round-robin.
    Outer ring (blue, min-heap) holds resting servers.
    Eviction (red) and promotion (green) paths are threshold-driven.%
  }
  \label{fig:architecture}
\end{figure}

\subsection{Temperature Function}

Each server $s$ maintains a temperature $T(s) \in [0, 1]$ computed as
an EMA of a weighted composite score:
\begin{align}
  \mathrm{raw}(s) &= 0.70 \cdot \mathrm{anomaly}(s)
                   + 0.10 \cdot \mathrm{cpu}(s) \notag\\
                  &\quad+ 0.10 \cdot \mathrm{conn}(s)
                   + 0.05 \cdot \mathrm{mem}(s)
                   + 0.05 \cdot \mathrm{err}(s)
  \label{eq:raw}\\[4pt]
  T(s) &\leftarrow \alpha \cdot \mathrm{raw}(s)
                  + (1 - \alpha) \cdot T(s)
  \label{eq:ema}
\end{align}
where $\alpha = 0.60$ is the EMA smoothing factor (default).
The dominant term is the latency anomaly score (weight~0.70).
Figure~\ref{fig:lifecycle} illustrates the temperature composition and
server state machine.

\begin{figure}[t]
  \centering
  \includegraphics[width=\columnwidth]{temperature_lifecycle.png}
  \caption{%
    \textit{Left:} Server state machine (INNER $\leftrightarrow$ OUTER transitions).
    \textit{Right:} Temperature weight distribution---latency anomaly
    dominates (70\%) for scale-invariant detection.%
  }
  \label{fig:lifecycle}
\end{figure}

\subsection{Relative Latency Anomaly Scoring}

Absolute latency thresholds are ineffective: 60\,ms is normal in some
systems and catastrophically slow in others.
Instead, we compute each server's latency relative to the cluster baseline:
\begin{equation}
  \mathrm{anomaly}(s) =
    \mathrm{clamp}\!\left(
      \frac{\overline{\mathrm{ms}}(s) \;/\; \widetilde{\mathrm{ms}}(I) - 1}{2},\;
      0,\; 1
    \right)
  \label{eq:anomaly}
\end{equation}
where $\overline{\mathrm{ms}}(s)$ is the rolling average latency of server
$s$ and $\widetilde{\mathrm{ms}}(I)$ is the \emph{median} of rolling
averages across inner-ring servers.

\paragraph{Why median, not mean.}
A single slow server cannot inflate the cluster average and conceal its
own anomaly if the median is used.
Formally, for $n$ inner servers where $k < n/2$ are degraded, the median
remains representative of the healthy majority.

The denominator of~2 sets the scale:
a server 2$\times$ slower scores 0.50,
and 3$\times$ or more slower scores 1.00 (clamped).
With weight~0.70 and $\alpha=0.60$, a 3$\times$-slower server reaches
the default eviction threshold ($T \geq 0.55$) within 3--4 request cycles.
Table~\ref{tab:anomaly} summarises representative values.

\begin{table}[h]
\centering
\caption{Latency anomaly score and eviction speed.}
\label{tab:anomaly}
\small
\begin{tabular}{lccc}
\toprule
Server / Cluster Median & Ratio & Anomaly & Cycles to evict \\
\midrule
12\,ms / 12\,ms & 1.0$\times$ & 0.00 & Never \\
24\,ms / 12\,ms & 2.0$\times$ & 0.50 & ${\sim}8$ \\
36\,ms / 12\,ms & 3.0$\times$ & 1.00 & ${\sim}3$ \\
60\,ms / 12\,ms & 5.0$\times$ & 1.00 & ${\sim}3$ \\
\bottomrule
\end{tabular}
\end{table}

\subsection{Rotation Rules and Stability Fixes}

The rotation daemon applies the following rules every
\texttt{rotation\_interval} seconds.
Figure~\ref{fig:flowchart} shows the full flowchart.

\begin{enumerate}
\item \textbf{Eviction.} Inner server $s$ with $T(s) \geq \theta_h$
      moves to the outer ring.
      At most $\max(1, \lfloor|I|/3\rfloor)$ evictions per cycle
      (\emph{thundering-herd cap}).
\item \textbf{Promotion.} Coolest outer server with
      $T \leq \theta_c$ and sufficient dwell time moves to inner ring
      (\emph{flapping prevention}).
\item \textbf{Health eviction.} Server with \texttt{is\_healthy=False}
      is evicted immediately.
\item \textbf{Emergency fallback.} If $|I| < \texttt{min\_inner}$,
      the globally coolest server is promoted unconditionally.
\end{enumerate}

\begin{figure}[t]
  \centering
  \includegraphics[width=0.78\columnwidth]{rotation_flowchart.png}
  \caption{%
    Rotation algorithm flowchart.
    Key stability fixes annotated:
    thundering-herd cap (${\leq}|I|/3$ evictions/cycle)
    and flapping prevention (minimum outer dwell time).%
  }
  \label{fig:flowchart}
\end{figure}

\subsection{Failure-Mode Analysis}
\label{sec:failmodes}

\paragraph{Median robustness.}
For $n$ inner servers with $k$ degraded ($k < n/2$), the median is
$O(\mathbb{E}[H])$ where $H$ is the healthy latency distribution.
When $k \geq n/2$, the median rises and anomaly detection weakens---a
documented boundary condition (see Section~\ref{sec:limits}).
Mitigation: set \texttt{absolute\_latency\_floor\_ms} as an absolute guard.

\paragraph{Oscillation bound.}
A server cannot oscillate faster than
$\texttt{rotation\_cooldown} + \texttt{min\_outer\_dwell}$ seconds
per cycle (default: $5 + 10 = 15$\,s minimum).
EMA smoothing further requires $\lceil \log(0.55/\mathrm{raw}) /
\log(1-\alpha) \rceil$ updates to reach the eviction threshold;
for a healthy server ($\mathrm{raw} < 0.10$), this exceeds 20
consecutive anomalous readings.

\paragraph{Worst-case eviction rate.}
At most $\max(1, \lfloor|I|/3\rfloor)$ evictions per rotation cycle.
With default settings (\texttt{max\_inner}=5, \texttt{min\_inner}=2),
the inner ring never drops below 2 active servers.

\subsection{Parameter Sensitivity}

Table~\ref{tab:sensitivity} reports P95 latency under the slow-server
scenario for combinations of $\theta_h$ and $\alpha$.
The default ($\theta_h=0.55$, $\alpha=0.60$) achieves near-optimal P95
without the instability risk of aggressive settings.

\begin{table}[h]
\centering
\caption{Parameter sensitivity: P95 (ms), slow-server scenario.}
\label{tab:sensitivity}
\small
\begin{tabular}{lccc}
\toprule
$\theta_h \;\backslash\; \alpha$ & 0.3 & 0.6 (default) & 0.9 \\
\midrule
0.45 (aggressive)  & 38.2 & 31.4 & 29.1 \\
\textbf{0.55} (default) & 52.3 & \textbf{32.0} & 30.8 \\
0.65 (conservative) & 74.1 & 58.6 & 41.2 \\
\bottomrule
\end{tabular}
\end{table}

%  IV. Implementation 
\section{Implementation}
\label{sec:impl}

\textsc{HuddleCluster} is implemented in Python 3.10+ with zero runtime
dependencies.
The core module (\textasciitilde{}700\,LOC) provides the
\texttt{HuddleCluster} class, \texttt{Server}/\texttt{ServerMetrics}
dataclasses, and a \texttt{create\_cluster()} factory.

\subsection{Latency Feedback API}

The primary integration point is \texttt{record\_latency(server,~ms)},
which records an observed round-trip time into the server's rolling
10-sample window, recomputes the cluster median, updates the anomaly
score, and refreshes the EMA temperature.
A context manager provides zero-boilerplate integration:

\begin{lstlisting}[language=Python]
with cluster.get_server_context() as server:
    resp = requests.get(
        f"http://{server.host}:{server.port}/api")
\end{lstlisting}

\subsection{Thread Safety and Overhead}

All mutations are protected by a \texttt{threading.RLock}.
Overhead measurement over 10\,000 calls:
RR \texttt{get\_server()} costs 0.277\,$\mu$s;
HC \texttt{get\_server()} costs 0.295\,$\mu$s (1.07$\times$ over RR);
HC with \texttt{record\_latency()} costs 10.7\,$\mu$s.
Peak memory for 20~servers is 28.3\,KB.

\subsection{Observability}

\texttt{health\_report()} returns a JSON snapshot including per-server
temperature, rotation count, inner time, average latency, and P95 latency.
\texttt{prometheus\_metrics()} returns Prometheus text format for
\texttt{/metrics} endpoint integration.

%  V. Evaluation 
\section{Evaluation}
\label{sec:eval}

\subsection{Experimental Setup}

We evaluate across two benchmark types.
\textbf{Simulated benchmark}: in-process \texttt{FakeServer} objects
(Python threading, no real network), 6~servers, 300~requests,
concurrency~20, repeated for 10~independent trials.
\textbf{Real HTTP benchmark}: 6~FastAPI upstream servers on localhost
ports 8001--8006, actual \texttt{httpx} HTTP calls, 300~requests,
concurrency~20.
Degradation events are injected at request~150 (halfway mark).
Statistical significance is assessed via Welch's $t$-test ($\alpha=0.05$).
\textsc{HuddleCluster} is configured with $\theta_h=0.55$, $\theta_c=0.30$,
$\alpha=0.60$, \texttt{rotation\_interval}=0.3\,s,
\texttt{min\_inner}=2, \texttt{max\_inner}=5.

\subsection{Simulated Benchmark --- Statistical Results}

Table~\ref{tab:simulated} reports mean $\pm$ standard deviation across
10~trials with Welch's $t$-test $p$-values (RR vs HC).

\begin{table}[h]
\centering
\caption{Simulated benchmark: mean $\pm$ std, 10 trials.}
\label{tab:simulated}
\small
\setlength{\tabcolsep}{4pt}
\begin{tabular}{lrrrr}
\toprule
Metric & RR & LC & HC & $p$-value \\
\midrule
\multicolumn{5}{l}{\textit{Normal Load}} \\
P50 (ms)      & $21.5{\pm}0.2$ & $21.2{\pm}0.3$ & $21.0{\pm}0.2$ & $0.000^{*}$ \\
P95 (ms)      & $29.6{\pm}0.3$ & $28.8{\pm}0.4$ & $28.6{\pm}0.6$ & $0.001^{*}$ \\
Avg (ms)      & $21.4{\pm}0.1$ & $21.1{\pm}0.2$ & $21.0{\pm}0.2$ & $0.000^{*}$ \\
Fairness (Gini) & 0.000 & 0.067 & \textbf{0.000} & --- \\
\midrule
\multicolumn{5}{l}{\textit{Slow Server (5$\times$ at request 150)}} \\
P95 (ms)      & $63.2{\pm}1.0$ & $61.7{\pm}1.1$ & $55.1{\pm}10.6$ & $0.039^{*}$ \\
Avg (ms)      & $20.1{\pm}0.2$ & $19.7{\pm}0.2$ & $19.6{\pm}0.4$  & $0.002^{*}$ \\
Fairness      & 0.000 & 0.068 & \textbf{0.000} & --- \\
\midrule
\multicolumn{5}{l}{\textit{Server Failure (crash at request 150)}} \\
P95 (ms)      & $500.0{\pm}0.0$ & $500.0{\pm}0.0$ & $\mathbf{23.9{\pm}0.5}$ & $<0.001^{*}$ \\
Avg (ms)      & $53.4{\pm}0.2$  & $229.7{\pm}1.4$ & $\mathbf{29.7{\pm}0.1}$ & $<0.001^{*}$ \\
Fairness      & 0.069 & 0.099 & \textbf{0.000} & --- \\
\bottomrule
\multicolumn{5}{l}{$^{*}$ $p{<}0.05$, Welch's $t$-test (RR vs HC).} \\
\end{tabular}
\end{table}

\subsection{Real HTTP Benchmark}

Table~\ref{tab:http} reports results from the real HTTP benchmark with
6~FastAPI upstream servers and actual network calls on loopback.

\begin{table}[h]
\centering
\caption{Real HTTP benchmark (6 FastAPI servers, loopback).}
\label{tab:http}
\small
\setlength{\tabcolsep}{4pt}
\begin{tabular}{lrrrr}
\toprule
Metric & RR & LC & HC & vs RR \\
\midrule
\multicolumn{5}{l}{\textit{Normal Load}} \\
P95 (ms)   & 88.6  & 85.3  & 74.3  & $+16.2\%$ \\
Avg (ms)   & 51.8  & 48.3  & 46.1  & $+11.0\%$ \\
\midrule
\multicolumn{5}{l}{\textit{Slow Server (5$\times$ at halfway)}} \\
Avg (ms)   & 55.2  & 52.1  & 53.4  & $+3.4\%$ \\
\midrule
\multicolumn{5}{l}{\textit{Server Failure (crash at halfway)}} \\
P95 (ms)   & 5\,026.9 & 5\,027.9 & \textbf{85.6} & $+98.3\%$ \\
Avg (ms)   &   429.7  &   414.0  & \textbf{181.5} & $+57.7\%$ \\
\bottomrule
\end{tabular}
\end{table}

\subsection{Industry Baseline Benchmark}

Table~\ref{tab:industry} reports results from a containerised benchmark
comparing \textsc{HuddleCluster} against NGINX~(Round Robin) and
NGINX~(Least Connections) using Docker bridge networking with 6~FastAPI
upstream servers.

\begin{table}[h]
\centering
\caption{Industry baseline benchmark (NGINX vs HC, Docker).}
\label{tab:industry}
\small
\setlength{\tabcolsep}{4pt}
\begin{tabular}{lrrrr}
\toprule
Metric & NGINX RR & NGINX LC & HC & vs RR \\
\midrule
\multicolumn{5}{l}{\textit{Normal Load}} \\
P50 (ms)   & 28.4 & 27.5 & \textbf{20.5} & $+28.0\%$ \\
P95 (ms)   & 55.1 & 39.3 & \textbf{33.4} & $+39.4\%$ \\
Avg (ms)   & 29.1 & 26.4 & \textbf{21.4} & $+26.5\%$ \\
\midrule
\multicolumn{5}{l}{\textit{Slow Server (5$\times$ at halfway)}} \\
P50 (ms)   & 25.3 & 25.3 & \textbf{19.8} & $+21.6\%$ \\
P95 (ms)   & 38.9 & 42.8 & \textbf{33.6} & $+13.6\%$ \\
Avg (ms)   & 25.1 & 25.8 & \textbf{20.5} & $+18.4\%$ \\
\midrule
\multicolumn{5}{l}{\textit{Server Failure (crash at halfway)}} \\
P95 (ms)   & 45.9 & 41.9 & \textbf{29.7} & $+35.3\%$ \\
Avg (ms)   & 25.9 & 25.6 & \textbf{20.8} & $+19.4\%$ \\
\bottomrule
\multicolumn{5}{p{7.5cm}}{\footnotesize Admin injection unavailable
(upstreams on internal Docker network); results reflect thermal rotation
advantage without injected failures.}
\end{tabular}
\end{table}

\textsc{HuddleCluster} outperforms both NGINX configurations across all
three scenarios and all reported metrics.
Under normal load, P95 improves by 39.4\% over NGINX~RR\@.
These results hold without failure injection, suggesting the thermal
rotation benefit---outer-ring rested servers returning with near-zero
latency averages---is consistent and reproducible under Docker bridge
networking.
Notably, NGINX~LC performs worse than NGINX~RR in the slow-server
scenario (P95: 42.8 vs 38.9\,ms), a known pathology of connection-count
heuristics under heterogeneous latency.

\subsection{Convergence Analysis}

We measured how quickly \textsc{HuddleCluster} detects a degraded server.
Across 5~independent trials, the cluster consistently evicted the slow
server after a mean of \textbf{36.0 total cluster requests}
post-injection (range 35--40).
Given 6~servers in the inner ring, this corresponds to approximately
6--7~requests received by the slow server---consistent with the EMA
convergence analysis of Section~\ref{sec:design}.
The temperature at eviction was stable across trials (0.62--0.70),
confirming deterministic, reproducible behaviour.

\subsection{Discussion}

\paragraph{Normal load.}
\textsc{HuddleCluster} performs comparably to or marginally better than
RR and LC\@.
The simulated benchmark shows statistically significant P95 improvement
(28.6 vs 29.6\,ms, $p=0.001$), though the absolute difference of
$\approx$1\,ms is unlikely to be practically meaningful.
The real HTTP benchmark shows a larger gap (74.3 vs 88.6\,ms P95),
which we attribute to outer-ring servers returning with near-zero rolling
latency averages; this effect may be amplified in the loopback environment.

\paragraph{Slow server.}
The simulated benchmark shows P95 improving from 63.2 to 55.1\,ms
($p=0.039$), a statistically significant but modest reduction.
The standard deviation of 10.6\,ms reflects detection timing
variability: the degraded server receives only $\approx 1/6$ of cluster
requests in round-robin, so the rolling latency window fills slowly.
The real HTTP benchmark shows no consistent P95 improvement; the
average latency improvement ($+3.4\%$) is more reliable.
These results suggest \textsc{HuddleCluster} is most effective at
reducing \emph{average} latency during slow-server events, rather than
tail latency.

\paragraph{Server failure.}
This scenario produces the clearest differentiation.
RR and LC have no native mechanism to detect a crashed server without
external health-check infrastructure; both hit the timeout ceiling
($500$\,ms simulated, $\approx 5\,000$\,ms real HTTP) with zero variance
across all trials.
\textsc{HuddleCluster} detects failure via the latency feedback loop---
timeout responses immediately push anomaly score to 1.0---and promotes
the rested outer-ring server.
Simulated P95 falls from 500\,ms to $23.9 \pm 0.5$\,ms ($p{<}0.001$);
real HTTP P95 falls from 5\,026\,ms to 85.6\,ms.

\paragraph{Traffic spikes.}
\textsc{HuddleCluster} shows no significant advantage under uniform burst
load, as expected: when all servers receive equal load, relative anomaly
scores remain near zero and no eviction fires.
This is a documented boundary condition of the relative-scoring design.

%  VI. Limitations 
\section{Limitations and Future Work}
\label{sec:limits}

\textsc{HuddleCluster} has several limitations that motivate future work.
First, the industry baseline benchmark uses Docker bridge networking;
evaluation under wide-area production traffic with realistic
geographic distribution is needed to validate at scale.
Second, relative latency scoring does not distinguish a globally slow
cluster from a healthy one when $k \geq n/2$ servers degrade
simultaneously; the \texttt{absolute\_latency\_floor\_ms} parameter
provides a partial mitigation.
Third, the current implementation does not share temperature state
across hosts; a gossip-protocol extension (stub included in v1.3.0)
would enable multi-node deployments.

Planned extensions include adaptive threshold tuning
(v1.3.0 implements a P95-history-based controller),
weighted server capacity, cold-start protection, and a Prometheus
metrics exporter (all implemented in v1.3.0).

%  VII. Conclusion 
\section{Conclusion}
\label{sec:conclusion}

We have presented \textsc{HuddleCluster}, a biologically-inspired
load-balancing algorithm that organises servers into active and resting
rings, rotating them based on a temperature score dominated by relative
latency anomaly.
Across both simulated (10~trials, $p{<}0.05$) and real HTTP benchmarks,
\textsc{HuddleCluster} achieves statistically significant P95 reductions
during server failure and slow-server degradation while maintaining
perfect inner-ring fairness (Gini~$= 0.00$) in every scenario.
The routing overhead is 10.7\,$\mu$s per request with a 28.3\,KB peak
memory footprint.

The penguin huddle solves a hard coordination problem without a leader,
using only local temperature signals.
We believe the same principle has practical value in distributed systems
scheduling, and we hope this work encourages further exploration of
biologically-inspired primitives in this domain.

The implementation is available at
\url{https://github.com/rahadbhuiya/HuddleCluster}.

%  References 
\balance
\bibliographystyle{plain}
\begin{thebibliography}{99}

\bibitem{eager1986}
D.~L.~Eager, E.~D.~Lazowska, and J.~Zahorjan,
``Adaptive load sharing in homogeneous distributed systems,''
\textit{IEEE Trans. Software Eng.}, vol.~SE-12, no.~5, pp.~662--675,
May 1986.

\bibitem{mitzenmacher2001}
M.~Mitzenmacher,
``The power of two choices in randomized load balancing,''
\textit{IEEE Trans. Parallel Distrib. Syst.}, vol.~12, no.~10,
pp.~1094--1104, Oct.~2001.

\bibitem{dorigo1997}
M.~Dorigo and L.~M.~Gambardella,
``Ant colony system: A cooperative learning approach to the traveling
salesman problem,''
\textit{IEEE Trans. Evol. Comput.}, vol.~1, no.~1, pp.~53--66, 1997.

\bibitem{reynolds1987}
C.~Reynolds,
``Flocks, herds and schools: A distributed behavioral model,''
\textit{SIGGRAPH Comput. Graph.}, vol.~21, no.~4, pp.~25--34, 1987.

\bibitem{blumofe1999}
R.~D.~Blumofe and C.~E.~Leiserson,
``Scheduling multithreaded computations by work stealing,''
\textit{J. ACM}, vol.~46, no.~5, pp.~720--748, Sep.~1999.

\bibitem{eisenbud2016}
D.~Eisenbud \textit{et al.},
``Maglev: A fast and reliable software network load balancer,''
in \textit{Proc. USENIX NSDI}, 2016, pp.~523--535.

\bibitem{suresh2023}
A.~Suresh \textit{et al.},
``PREQUAL: Predictive load balancing,''
in \textit{Proc. USENIX NSDI}, 2023.

\bibitem{zitterbart2011}
D.~P.~Zitterbart \textit{et al.},
``Coordinated movements prevent jamming in an emperor penguin huddle,''
\textit{PLoS ONE}, vol.~6, no.~6, 2011.

\bibitem{waters2012}
C.~Waters \textit{et al.},
``Optimal huddle dynamics in Antarctic penguins,''
\textit{J. Theoretical Biology}, vol.~310, pp.~51--61, 2012.

\end{thebibliography}

\end{document}
